\clearpage
\VerdierPhysicalPageAnchor{0184}%
\begingroup\small\setlength{\parskip}{0pt}\setlength{\emergencystretch}{1.5em}
\setlength{\abovedisplayskip}{4pt plus 2pt minus 2pt}%
\setlength{\belowdisplayskip}{4pt plus 2pt minus 2pt}%

\VerdierSemanticAnchor{ch3d-1-3-4-c}
\noindent\textit{c) $D(\mathcal A)$ 的每个特异三角都同构于形如
$(X,Y,Z,[u],[v],\delta(S))$ 的特异三角，其中
$S=0\longrightarrow X\xrightarrow{u}Y\xrightarrow{v}Z\longrightarrow0$
是 $\mathcal A$ 中的一条正合列。}

\VerdierSemanticAnchor{ch3d-1-3-4-proof}
a) 和 b) 的证明留给读者。下面证明 c)。由（第二章，2.2.6），
$D(\mathcal A)$ 的每个特异三角都同构于 $K(\mathcal A)$ 的一个特异三角在
$Q\colon K(\mathcal A)\longrightarrow D(\mathcal A)$ 下的像。由
（第一章，3.3.1）和公理 TRII，$K(\mathcal A)$ 的每个特异三角都同构于
一个如下形式的三角：

\VerdierSemanticAnchor{ch3d-1-3-4-cone-triangle}
\[
Y\xrightarrow{\dot{p}(u)}c(\dot u)
\xrightarrow{\dot{q}(u)}X[1]
\xrightarrow{-\dot u[1]}Y[1].
\]
因此，$D(\mathcal A)$ 的每个特异三角都同构于一个如下形式的三角：

\VerdierSemanticAnchor{ch3d-1-3-4-derived-triangle}
\[
(Y,c(\dot u),X[1],[p(u)],[q(u)],-[u][1]).
\]
现注意到：

\VerdierSemanticAnchor{ch3d-1-3-4-cone-exact}
\[
S=0\longrightarrow Y\xrightarrow{p(u)}c(u)
\xrightarrow{q(u)}X[1]\longrightarrow0
\]
是 $\operatorname{comp}(\mathcal A)$ 中的一条正合列。于是由
（第二章，1.2.4），存在从三角
\[
(Y,c(\dot u),X[1],[p(u)],[q(u)],-[u][1])
\]
到三角
\[
(Y,c(\dot u),X[1],[p(u)],[q(u)],\delta(S)).
\]
的一个同构。

\VerdierSemanticAnchor{ch3d-1-3-5}
\medskip\noindent\textbf{命题 1.3.5.}\quad
\textit{a) 每个函子
$F\colon\operatorname{comp}(\mathcal A)\longrightarrow C$，若把拟同构变为
$C$ 中的同构，就唯一地通过函子
$\operatorname{comp}(\mathcal A)\longrightarrow D(\mathcal A)$ 分解。}

\noindent\textit{b) 从 $\operatorname{comp}(\mathcal A)$ 到阿贝尔范畴
$B$ 的每个半正合函子，若把拟同构变为同构，就唯一地分解为典范
$\delta$-函子
$\operatorname{comp}(\mathcal A)\longrightarrow D(\mathcal A)$ 与一个从
$D(\mathcal A)$ 到 $B$ 的上同调函子的复合。}

\VerdierSemanticAnchor{ch3d-1-3-5-proof}
断言 b) 由断言 a) 以及 (1.3.4) 给出的 $D(\mathcal A)$ 中特异三角的描述
推出。断言 a) 由 $D(\mathcal A)$ 的定义（第二章，2.2.6）和下面的引理推出。

\VerdierSemanticAnchor{ch3d-1-3-6}
\medskip\noindent\textbf{引理 1.3.6.}\quad
\textit{设 $\mathcal A$ 是一个加法范畴（不必是阿贝尔范畴），$F$ 是从
$\operatorname{comp}(\mathcal A)$ 到范畴 $C$ 的一个函子。下列条件等价：}
\endgroup
